\section{Completion of the actual parameter record}
\label{sec:parameter-completion}

\status{PN2--PN15 prove the fourteen parameter families retained after RF16.
Together with RF12--RF14 and RF16, PN16 constructs every field of the
literal \texttt{ParameterChoice'} record, with its prescribed constant
list, under the supplied \texttt{Val} and \texttt{HostArith} hypotheses.
The threshold is uniform in the actual datum and size datum. This closes
\texttt{explicitConstants\_parameterConditions\_large}; the construction
of the host, columns and initial size datum remains a separate input.
The \texttt{SizesAt} and \texttt{NumAt} joins are treated in the following
component.}

\paragraph{PN1: common data, exact scales, and the order of thresholds.}
Fix $0<\epsilon<1/52$ and a standard constant record $C$. All assertions
in this section have the source's \texttt{SzRow} order
\[
 \exists c_{D,0}\ \forall c_D\geq c_{D,0}\
 \exists c_{h,0}\ \forall c_h\geq c_{h,0}\
 \exists x_0\ \forall D,L:\quad
 \left.
 \begin{array}{c}
 L.\mathrm{sc}.\epsilon=\epsilon,\quad L.\mathrm{sc}.x\geq x_0,\\
 \operatorname{Val}(C,c_D,c_h,L)
 \end{array}\right\}\Longrightarrow\hbox{the stated conclusion}.
                                                               \tag{PN.1}
\]
Here $L$ is an actual \texttt{SizeData\textsubscript{2}} over the same
actual $D$. Its windows and record identities are available; no new
geometric object is chosen in a numerical proof. The eventual threshold
$x_0$ may depend on $\epsilon,C,c_D,c_h$, but not on $D,L$ or on the
auxiliary scale parameter $J$.

Use the notation
\begin{gather*}
 n=D.R.n,\quad f=2n,\quad r=k+m-1,\quad
 N=N_m=k+m(t+1),\quad \mathcal L=x^\epsilon,\quad\ell=\log x,\\
 b=e^{\mathcal L},\quad a=xb,\quad X'=e^{x^{\epsilon/3}},\quad
 b'=\frac{xb}{X'},\quad w=\frac{b}{nX'},\quad X''=xnX'=\frac{a}{w},\\
 \Delta_{\rm cl}=\lceil x^{c_D}\rceil_+,\quad
 H=2\lceil x^{c_h}\rceil_+,\quad
 S=4k\Delta_{\rm cl},\quad\eta=mH,\quad h_0=H,\quad
 d_2=\frac{|T_2|}{n},\quad\lambda=L.\mathrm{LaG}.
                                                               \tag{PN.2}
\end{gather*}
Thus $S=L.\mathrm{Sc}$ and $\eta=L.\mathrm{etaSz}$, by the actual
\texttt{Sc\_eq}, \texttt{Del\_eq} and \texttt{eta\_eq} identities.
The scalar $\Delta_{\rm cl}$ is unrelated to an identification of the
integral Weyl operator $\Delta$ with its dual $\nabla$.
Set
\begin{align*}
 A_*&=nt^2(N+t)\mathcal L^2+kn^2(S+t+\mathcal L),\\
 \ell_*&=\frac{mb}{A_*},\qquad
 \ell_A=\frac{k\mathcal L}{md_2},\qquad
 r_S=\frac m2-1,\qquad M=r+mt=N-1.
                                                               \tag{PN.3}
\end{align*}
These are the literal \texttt{lStar}, \texttt{lA}, \texttt{rS} and
\texttt{Mrt}; $m$ is even and at least two. Also
$\texttt{qSig}=m$. All denominators in (PN.3) are positive.

RF12 gives, after a threshold depending only on $\epsilon$,
\[
 1\leq n\leq f\leq x^3,\quad
 2\leq m\leq x,\quad x^{1/4}\leq m,\quad
 k\leq r\leq x^5,\quad
 1\leq t\leq x^{401},\quad N\leq x^{402},\quad
 1\leq\mathcal L\leq x.                                    \tag{PN.4}
\]
The weaker lower bound $m\geq x^{1/4}$ follows from RF12's exact
$m\geq x^{1/2-\epsilon/16}$. The supplied second-set bound gives
$d_2\geq x^{-\epsilon}\geq x^{-1}$, while $d_2\leq1$ is already a
size-datum field. Consequently $\ell_A\leq x^7$. RF14 gives
$\mathcal L\leq\lambda\leq B\mathcal L\ell^2$ for a fixed $B>0$;
RF11 therefore gives $1\leq\lambda\leq x$ eventually.
For $c_D,c_h\geq0$ and $x\geq8$, rounding gives
\[
 x^{c_D}\leq\Delta_{\rm cl}\leq2x^{c_D},\qquad
 x^{c_h}\leq H\leq4x^{c_h},\qquad
 S\leq x^{c_D+6},\qquad \eta\leq x^{c_h+3}.             \tag{PN.5}
\]

We repeatedly use the already proved RF11 comparison: for every fixed
real exponent $v$, every positive fixed coefficient, and
$1\leq n\leq x^3$, a sufficiently large common threshold makes any
prescribed polynomial in $x$ smaller than $b,b',w$; it also gives
$b'x^v\leq b$ and $x^vX'\leq b$. The coefficient is absorbed using a
positive exponent margin. These are uniform statements in the actual
$n,D,L$, not estimates obtained by specializing to one datum.

\paragraph{PN2: both host primes and the exact conductor-product bound.}
This proves family [1]. Put $C_8=C.P.C_{\rm pde}\geq1$,
$\zeta=\epsilon/J$, and $Q_c=QS\,q\,d=\operatorname{condE}(D)$.
The supplied \texttt{HostVal} gives $1\leq J\leq C_8$ and a real $Y$
with $x^{1/J}\leq Y\leq x$. Its \texttt{PairRow} clauses say
\begin{gather*}
 e^{Y^{\epsilon/5}/2}<q\leq e^{Y^{\epsilon/5}},\qquad
 e^{Y^{\epsilon/5}(\log Y)^2/2}<d
                 \leq e^{Y^{\epsilon/5}(\log Y)^2},\\
 \log Q_c\leq2Y^{\epsilon/5}(\log Y)^2.
\end{gather*}
Taking logarithms, and using $\ell/J\leq\log Y\leq\ell$, proves
\begin{gather*}
 \tfrac12x^{\zeta/5}<\log q\leq x^{\epsilon/5},\qquad
 \tfrac12x^{\zeta/5}(\ell/J)^2<\log d
                                  \leq x^{\epsilon/5}\ell^2,\\
 \log Q_c\leq2x^{\epsilon/5}\ell^2.                         \tag{PN.6}
\end{gather*}
For $\ell\geq C_8$, the lower bound for $\log d$ is at least the
lower bound displayed for $\log q$. Since
$\zeta\geq\epsilon/C_8$, RF11 makes
\[
 (2C.C_t+5)\ell<\tfrac12x^{\epsilon/(5C_8)}
                                      \leq\tfrac12x^{\zeta/5}.
\]
Thus $x^{2L.C_t+5}<q,d$, using $L.C_t=C.C_t=401$.
Requiring $x^{\epsilon/(30C_8)}\geq2$ also gives
\[
 x^{\zeta/6}\leq\tfrac12x^{\zeta/5}<\log q.            \tag{PN.7}
\]
All thresholds are uniform in $J$ because its upper bound is fixed.

By (PN.4),
\[
 \log(Q_c^{36}t^{150})
 \leq72x^{\epsilon/5}\ell^2+60150\ell
 \leq\tfrac12x^\epsilon                                  \tag{PN.8}
\]
eventually, by the positive exponent gaps in RF11. Exponentiating
gives $Q_c^{36}t^{150}\leq b^{1/2}$. The host inequalities
$p_{\rm de},p'_{\rm de}\leq Q_c^{C_8}$ similarly give
\[
 \log p_{\rm de},\ \log p'_{\rm de}
       \leq2C_8x^{\epsilon/5}\ell^2\leq\mathcal L.       \tag{PN.9}
\]
Finally $|L.\mathrm{logDerivativeAtOne}|\leq x^\epsilon$ is precisely
the supplied \texttt{Val.dlog} field. This completes the five-part
family. Here \texttt{condE} is the defined product $QSqd$; the proof
does not construct a host field or prove an analytic logarithmic-
derivative estimate from the numerical inequalities.

\paragraph{PN3: the four logarithmic and block-window inequalities.}
This proves family [4]. For $m\geq4$, evenness gives
$r_S=m/2-1\geq m/4$. RF11 and (PN.4) imply, eventually,
\[
 0\leq\log f\leq3\ell,\qquad
 9\ell^2\leq\tfrac14x^{1/4}\leq\tfrac m4\leq r_S.
                                                               \tag{PN.10}
\]
Hence $\log f\leq\sqrt{r_S}$ and
$(\log f)^2\leq r_S\leq m\leq m^2\mathcal L$.
The standard constants give $c_{LL}=C.P.c_{LL}>0$.
Since $r\geq k\geq x$, requiring $x\geq c_{LL}^{-1}$ gives
$c_{LL}r\geq1$ and $\log r\geq\ell$. Therefore
\[
 \log\log b=\epsilon\ell\leq\ell
                    \leq c_{LL}rf\log r.                    \tag{PN.11}
\]
The actual block expression, with $C.P.C_{nW}=10$, is
\[
 n_W(L,10)=10r\left(f+\frac{f^2(k-1)^3\mathcal L^2}{m}\right)
             \geq10rf\geq(4f+10)r,                         \tag{PN.12}
\]
because $f\geq2$. These are all four fields of
\texttt{szRow\_logf}; no conductor-field identity is needed for the
source field named \texttt{log\_sq\_conductor\_le}.

\paragraph{PN4: all five bounds involving the height scale.}
This proves family [5]. Choose $c_h\geq401$. Then
$t\leq x^{401}\leq x^{c_h}\leq\eta$.
The upper bound in (PN.5) and the rank bounds give
\begin{align*}
 \eta t^3fk^3&\leq x^{c_h+3+1203+3+15}=x^{c_h+1224},\\
 xt\eta&\leq x^{c_h+405},\\
 m(t+1)\eta&\leq x^{c_h+406},
\end{align*}
where $t+1\leq x^{402}$ after $x\geq2$.
RF11 makes $x^{c_h+1224}\leq b,w$, and also
$(c_h+3)\ell\leq x^{\epsilon/4}$. Since $\eta>0$,
\[
 \boxed{\quad
 \eta\leq\frac{b}{t^3fk^3},\quad
 \log\eta\leq x^{\epsilon/4},\quad
 t\leq\eta,\quad t\eta\leq\frac wx,\quad
 m(t+1)\eta\leq w.\quad}                                  \tag{PN.13}
\]
The height exponent was fixed before these exponential comparisons
were imposed. Increasing the eventual $x$ threshold does not change it.

\paragraph{PN5: the four coupled height and cloud bounds.}
This proves family [6]. For $c_D,c_h\geq0$ put
$Q=x^{c_D+c_h+402}$ and take $x\geq8$.
Then $k,m,r,N,\mathcal L\leq Q$, $\Delta_{\rm cl}\leq2Q$,
$H\leq4Q$, and consequently
\[
 S\leq8Q^2\leq Q^3,\qquad
 \eta\leq4Q^2\leq Q^3,\qquad d_2Q\geq1.
\]
RF11 supplies $w\geq Q^{10}$. Since $m\geq1$,
\[
 d_2mw\geq Q^9,\qquad w\geq1.
\]
The four desired comparisons now follow by ordinary multiplication:
\begin{align*}
 m(S+\eta)&\leq2Q^4\leq Q^{10}\leq w\mathcal L,\\
 N\eta&\leq Q^4\leq Q^9\leq d_2mw,\\
 \eta^2r\mathcal L&\leq Q^8\leq Q^9\leq d_2mw^2,\\
 kS&\leq Q^4\leq w.
                                                               \tag{PN.14}
\end{align*}
The first three are also exactly the corresponding coupled
\texttt{SizeEta} clauses; the fourth is the parameter field
\texttt{kS\_le}.

\paragraph{PN6: the cloud upper bounds and its logarithm.}
This proves family [8]. Using $k\leq x^5$ and
$\Delta_{\rm cl}\leq2x^{c_D}$ gives
\[
 4kS\leq32x^{c_D+10},\qquad
 4C.P.C_{\Delta G}k^2\Delta_{\rm cl}
                       \leq8C.P.C_{\Delta G}x^{c_D+10}.
\]
Take $x\geq\max(32,8C.P.C_{\Delta G})$ and use RF11 to obtain
$x^{c_D+11}\leq b'$. Because $\lambda\geq1$, both right-hand
sides are at most $\lambda b'$. Dividing the first by the positive
$k$ yields
\[
 4S\leq\frac{\lambda b'}k,\qquad
 4C.P.C_{\Delta G}k^2\Delta_{\rm cl}\leq\lambda b'.       \tag{PN.15}
\]
Also
$\log\Delta_{\rm cl}\leq\log2+c_D\ell
\leq\tfrac12x^{\epsilon/4}$ eventually, by RF11.
This checks the literal half factor in \texttt{logDel}.

\paragraph{PN7: actual integer window length and logarithmic depth.}
This proves family [9]. The supplied identities and standard constants
are
\[
 L_w=10\lceil405\ell\rceil_++2\lceil30/\epsilon\rceil_++22,
 \quad C.P.C_{Lw}=4455,\quad C.P.C_{j0}=405,
 \quad j_0=\left\lceil\log_5\frac{nN}{m}\right\rceil_+.
                                                               \tag{PN.16}
\]
For every natural $l\geq11$, choose the actual integer
$j=\lceil7\ell\rceil_+$. Since $l-1\geq10>e$,
\[
 \ell_A\leq x^7=e^{7\ell}\leq e^j\leq(l-1)^j.
\]
Moreover $j\leq\lceil405\ell\rceil_+$, so the exact displayed
formula for $L_w$ gives $2j+2\leq L_w$. The ceiling inequalities give
\[
 L_w<4050\ell+60/\epsilon+34\leq4455\ell
\]
once $\ell\geq60/\epsilon+34$.

The rank inequality $k+m-1+5\leq t$ implies
$N=k+m(t+1)\leq2mt$. Thus $1<nN/m\leq2x^{404}$ by (PN.4).
Using $\log5>1$ and $\log2\leq\log5$, for $\ell\geq2$,
\[
 j_0<\frac{\log2+404\ell}{\log5}+1
         \leq404\ell+2\leq405\ell.                         \tag{PN.17}
\]
The elementary inequality $\log5>1$ follows, for example, from the
exponential series bound $e<3<5$. This proves the full integer-witness,
window-length and depth family, with all rounding preserved.

\paragraph{PN8: the complete exact column and arithmetic tuple.}
This proves family [10]. The supplied \texttt{HCVal} gives integers
$B_\tau$ with
\[
 v_\tau(i)=(k-2-i)4\Delta_{\rm cl},\quad
 b_\tau(j)=B_\tau+jH,\quad e_\tau(j)=b_\tau(j)+H/2.      \tag{PN.18}
\]
At a top corresponding to $t_T\in TT$,
$B_\tau=b_T(t_T)+\lfloor b'/x\rfloor$, which gives the exact
\texttt{topPlace\_first\_bound}. Off the top set, \texttt{HCVal}
also gives
\[
 u_\tau(0)=(k-2)4\Delta_{\rm cl}+k4\Delta_{\rm cl}+\tfrac12,
 \qquad u_\tau(k-1)=-\tfrac12.
\]
Consequently
\[
 u_\tau(0)-v_\tau(0)=S+\tfrac12\leq2S,
 \qquad u_\tau(k-1)=v_\tau(k-2)-\tfrac12.                 \tag{PN.19}
\]
Here $S\geq4$ since the natural ceiling $\Delta_{\rm cl}$ is at
least one. Adjacent $b$ coordinates differ by exactly $H=h_0$, and
every difference $e_\tau(j)-b_\tau(j)$ is exactly $H/2$.
Thus the adjacent-gap interval $[h_0,h_0+1]$ and both centering
errors, including the last coordinate, hold with error zero.

For each $j$, the last \texttt{HCVal} clause supplies
$\operatorname{MUQ}(\varphi_\tau,e_\tau(j))$ for the actual embeddings
$\varphi_\tau=D.\mathrm{host}.\mathrm{emb}(D.eTau(\tau))$.
By definition it means that there are integers $a_\tau$ with
$e_\tau(j)=a_\tau$ and, for every root of unity $\xi$ in the host field,
\[
                    \prod_\tau\varphi_\tau(\xi)^{a_\tau}=1.
                                                               \tag{PN.20}
\]
An integer considered as a rational has numerator that integer.
Replacing $a_\tau$ by $\operatorname{num}(e_\tau(j))$ therefore gives
exactly the host's \texttt{MuCond} field \texttt{e\_mu}; no parity
condition is silently substituted for this root-of-unity character
condition.

The remaining exact fields are as follows. The equality
$\eta_F=6x^{-C.P.C_F}$ is \texttt{Val.etaF\_eq}, and the standard
$C.P.C_{\eta F}=6$ gives its required coefficient. The two congruences
$p_{\rm de},p'_{\rm de}\equiv1\pmod{15}$ and their inequalities with
both $g_i$ are the supplied \texttt{Host.pde\_mod} and
\texttt{Host.pde\_ne\_g}. Finally \texttt{HostVal} says
$cc_K=\mathcal O_K$, so its absolute ideal norm is one:
\[
       \log\operatorname{Nm}(cc_K)=0\leq\mathcal L,
       \qquad \operatorname{Nm}(cc_K)=1\leq C.P.C_{ccK}\,d,
                                                               \tag{PN.21}
\]
using $C.P.C_{ccK}=1$ and the supplied positive prime $d$.
This supplies all twelve parts of \texttt{szRow\_exact}. It transports
the given character and prime conditions; it is not their existence proof.

\paragraph{PN9: every low node is quantitatively far from every midpoint.}
This proves family [11]. Choose $c_D\geq402$, $c_h\geq7c_D$,
and $Q=x^{c_D}\geq28$. From (PN.4), each of $n,t,N,\mathcal L,k$
is at most $Q$, and $\Delta_{\rm cl}\leq2Q$. Thus
\[
 A_*\leq2Q^6+8Q^5+2Q^4\leq12Q^6,
 \qquad 2A_*+t/2+3\leq28Q^6\leq Q^7.                   \tag{PN.22}
\]
Since $b/\ell_*=A_*/m\leq A_*$, and $H/2=\lceil x^{c_h}\rceil_+$
is at least $Q^7$,
\[
                 2b/\ell_*+t/2+3\leq H/2.
\]
For integer indices $i,j$, (PN.18) gives
\[
 |b_\tau(i)-e_\tau(j)|=|i-j-\tfrac12|H\geq H/2,
                                                               \tag{PN.23}
\]
because an integer has distance at least $1/2$ from $1/2$.
This proves the required inequality simultaneously for every
$\tau$ and every pair of allowed indices.

\paragraph{PN10: all cross-weight ratios and the final-bottom bound.}
This proves family [12]. RF11 lets us require
\[
               \eta+1\leq w\leq b'\leq b,
               \qquad 8x^6b'\leq b,\qquad x\geq8.        \tag{PN.24}
\]
Indeed $\eta\leq x^{c_h+3}$, while $w$ dominates every fixed
power; $w=b'/(xn)$; and $b'x^7\leq b$ suffices for the second
inequality. The supplied top upper bound, $\delta\leq1$,
$\mathcal L,\lambda\leq x$, $n\leq x^3\leq x^4$ and $k\geq1$
give
\[
              b_T(t_T)\leq b'+2x^6b'.                       \tag{PN.25}
\]
Apply the three actual \texttt{SizeData\textsubscript{2}} windows.
At a top their lower endpoint is at least $b'$, and their upper
endpoint is
\[
 b_T(t_T)+b'/x+\eta\leq3b'+2x^6b'\leq5x^6b'\leq2b.
\]
At a thick place the window is $[b,b+\eta+1]\subseteq[w,2b]$;
at a thin place it is $[w,w+\eta]\subseteq[w,2b]$.
The three cases exhaust the actual positions. Therefore
\[
          w\leq w_j(\tau)\leq2b
                  \quad(1\leq j\leq m).                   \tag{PN.26}
\]
Since $X''w=xb$ and $x\geq2$, this implies
$w_j(\tau)\leq X''w_{j'}(\tau')$ for every permitted pair.

For the second assertion of the family, \texttt{HCVal} gives
$B_\tau\geq w$ in each case: at a top,
$B_\tau=b_T+\lfloor b'/x\rfloor\geq b'$; at a thick place,
$B_\tau\geq\lceil b\rceil$; at a thin place,
$B_\tau=\lceil w\rceil$. Therefore
\[
 b_\tau(m-1)=B_\tau+(m-1)H
                \leq B_\tau+\eta\leq2B_\tau=2b_\tau(0).
                                                               \tag{PN.27}
\]
This proof uses the already supplied weight windows. It does not use
the desired parameter record to create an initial size datum.

\paragraph{PN11: all cloud-unit and count lower bounds.}
This proves family [14]. Write
$c=C.P.C_\Delta$, $F=C.P.C_F$ and $K=C.P.C_{\rm Count}$;
standardness gives $c,F,K\geq1$. Choose
\[
                 c_D\geq5c+Fc+406K+412,
                 \qquad x\geq\max(1,c).                    \tag{PN.28}
\]
The five literal left-hand sides are bounded as follows:
\[
\begin{array}{c|c}
 \text{left-hand side}&\text{upper bound}\\ \hline
 k^c(cm)^3/\eta_F^c&x^{5c+6+Fc}\\
 cmN&x^{404}\\
 k^c(cm^4\mathcal L)^4&x^{5c+24}\\
 Nm^{10}&x^{412}\\
 (Nf\mathcal L)^K&x^{406K}
\end{array}                                                   \tag{PN.29}
\]
For the first row, use $k^c\leq x^{5c}$,
$(cm)^3\leq x^6$ and
$\eta_F^c=(6x^{-F})^c\geq x^{-Fc}$.
For the third, $cm^4\mathcal L\leq x^6$.
The other three rows follow directly from (PN.4).
Every displayed exponent is at most the right side of (PN.28),
and hence every row is at most $x^{c_D}\leq\Delta_{\rm cl}$.
This proves all four conjuncts of \texttt{Del\_ge} and the separate
\texttt{Del\_count} field with their actual constants and real powers.

\paragraph{PN12: all ten node-height bounds, with a stronger first row.}
This proves family [15]. The following calculation also permits the
coefficients in the two marked rows to vary with $x$: assume
$0\leq A_r,A_k\leq x$, $x\geq200$, and put $Y=x^{c_D}\geq1$.
Choose
\[
                  c_h\geq\max(1300+c_D,406K),
                  \qquad K=C.P.C_{\rm Count}.                \tag{PN.30}
\]
Then $H\geq x^{1300}Y$. The two parts of $A_*$ satisfy
\[
 nt^2(N+t)\mathcal L^2\leq2x^{1209},\qquad
 kn^2(S+t+\mathcal L)\leq3x^{412}Y,
\]
because $S\leq x^6Y$ and $t,\mathcal L\leq x^{401}Y$.
In particular $A_*\leq5x^{1299}Y$, so
\[
 12A_*+2t+8\leq70x^{1299}Y\leq x^{1300}Y\leq H.
                                                               \tag{PN.31}
\]
This is stronger than the literal first row,
$12b/\ell_*+2t+8=12A_*/m+2t+8\leq H$.
The remaining nine left-hand sides have the explicit bounds
\[
\begin{array}{c|c}
 kt\mathcal L&x^{407}\\
 A_r r^{19}f^5\mathcal L&x^{112}\\
 (Nf\mathcal L)^K&x^{406K}\\
 4N^3+2t+6&12x^{1206}\leq x^{1207}\\
 A_k k m^{9/2}&x^{11}\\
 2M^3+t+3&6x^{1206}\leq x^{1207}\\
 120r\,q_{\rm sig}^{3/2}&120x^7\leq x^8\\
 Nm^{10}&x^{412}\\
 16q_{\rm sig}N&16x^{403}\leq x^{404}.
\end{array}                                                   \tag{PN.32}
\]
Here $M=N-1\leq N$, $q_{\rm sig}=m$,
$m^{9/2}\leq m^5$ and $m^{3/2}\leq m^2$ because $m\geq1$.
Every row is at most $H$ by (PN.30). To obtain the exact parameter
family, take $A_r=C.P.C_{h0r}$ and $A_k=C.P.C_{h0k}$, enlarging the
$x$ threshold to dominate these fixed constants. For the size-condition
application, the same calculation allows
$A_r=A_k=x^{\epsilon/100}\leq x$ without changing the exponent budget.
No inequality involving the desired parameter record has been used.

\paragraph{PN13: the additional cubic-rank and degree-power bounds.}
This proves family [16]. Let $K_N=C.P.C_{h0N}\geq1$ and choose
\[
                   c_h\geq\max(1210,3+402K_N).
                                                               \tag{PN.33}
\]
For $x\geq4$, the rank bounds imply
\[
 k+m(t+r+7)\leq x^5+3x^{402}\leq4x^{402}\leq x^{403}.
\]
It follows that
\[
 t+2\{k+m(t+(r+6)+1)\}^3+1
       \leq4x^{1209}\leq x^{1210}\leq H.                   \tag{PN.34}
\]
Also $fN^{K_N}\leq x^{3+402K_N}\leq H$.
These are the two parameter fields
\texttt{heightBound\_ge\_rank\_cubic} and
\texttt{heightBound\_ge\_degree\_power}; they are distinct from
RF15's four \texttt{ParamsT} height fields.

\paragraph{PN14: the full rank-bottom-cloud inequality.}
This proves family [18]. Use RF11 to require
\[
              b'x^{14}\leq b,
              \qquad x^{c_D+20}\leq b.
\]
Since $S\leq x^{c_D+6}$, both $b'$ and $S$ are at most $b/x^{14}$.
The uniform bounds $r\leq x^5$, $\ell\leq x$, $\ell_A\leq x^7$
give
\[
 r(b'+S)\ell\ell_A
       \leq x^{13}\frac{2b}{x^{14}}
       =\frac{2b}{x}\leq d_2mb,                           \tag{PN.35}
\]
where the last inequality uses $d_2\geq x^{-1}$ and $m\geq2$.
This proves the actual \texttt{rbS} field. The same inequality is the
last clause of \texttt{SizeEta}, since its \texttt{szlA} is $\ell_A$.

\paragraph{PN15: both strict lower-scale inequalities.}
This proves family [19]. Set
\[
 Q=x^{c_D+c_h+402},\qquad P=Q^7,
\]
and take $x\geq100$. Each of $n,t,N,\mathcal L,k,m$ is at most $Q$,
$\Delta_{\rm cl}\leq2Q$ and $H\leq4Q$. As in (PN.22),
\[
 0<A_*\leq2Q^6+8Q^5+2Q^4\leq P,\quad
 S\leq8Q^2\leq P,\quad\eta\leq4Q^2\leq P,\quad t\leq P.
\]
RF11 gives both $P^2\leq b'$ and $b'P^2\leq b$ at one threshold.
Consequently
\[
 \frac b{\ell_*}=\frac{A_*}{m}\leq P,
 \qquad \ell_*=\frac{mb}{A_*}\geq\frac bP
                  \geq b'P\geq100b'.                      \tag{PN.36}
\]
Because $P\leq b'$ and $b'>0$,
\begin{align*}
 b'+S+\eta+b/\ell_*&\leq b'+3P\leq4b'<100b'\leq\ell_*,\\
 2\eta+t+8b/\ell_*+7&\leq11P+7\leq12P
                                      \leq12b'<100b'\leq\ell_*.
                                                               \tag{PN.37}
\end{align*}
This preserves the strict signs in both source fields.
Since \texttt{szlStar} and \texttt{lStar} have the same definition,
these are also precisely the two \texttt{SizeLstar} inequalities.

\paragraph{PN16: construct all 111 fields of the compatible parameter record.}
Write $c=C.P.C_\Delta$, $F=C.P.C_F$, $K=C.P.C_{\rm Count}$ and
$K_N=C.P.C_{h0N}$. One sufficient common initial choice is
\begin{align*}
 c_{D,0}&=\max\{800,\ 5c+Fc+406K+412\},\\
 c_{h,0}(c_D)&=\max\{1000,\ 7c_D,\ 1300+c_D,\
                                 406K,\ 3+402K_N\}.
                                                               \tag{PN.38}
\end{align*}
For every $c_D\geq c_{D,0}$ and $c_h\geq c_{h,0}(c_D)$, take the
maximum of the finitely many eventual thresholds in PN1--PN15 and
RF12--RF14, RF16. This is allowed by (PN.1): every threshold precedes
the choice of $D,L$. The same chosen $D,L$ then satisfies all twenty-one
families. Their allocation is exact:
\[
\begin{array}{c|c@{\qquad}c|c}
 \text{family}&\text{proof}&\text{family}&\text{proof}\\ \hline
 1&\mathrm{PN2}&2,3&\mathrm{RF12}\\
 4&\mathrm{PN3}&5&\mathrm{PN4}\\
 6&\mathrm{PN5}&7&\mathrm{RF14}\\
 8&\mathrm{PN6}&9&\mathrm{PN7}\\
 10&\mathrm{PN8}&11&\mathrm{PN9}\\
 12&\mathrm{PN10}&13&\mathrm{RF13}\\
 14&\mathrm{PN11}&15&\mathrm{PN12}\\
 16&\mathrm{PN13}&17&\mathrm{RF13}\\
 18&\mathrm{PN14}&19&\mathrm{PN15}\\
 20&\mathrm{RF16}&21&\mathrm{RF14}.
\end{array}                                                   \tag{PN.39}
\]

Copy each of the following 23 real constants from $C.P$ into $P$:
\begin{gather*}
 C_{kLow},C_{k6},C_{fL},c_{LL},C_{nW},C_\rho,C_{\Delta G},
 C_{Lw},C_{j0},C_w,c_w,C_\Delta,\\
 C_{\rm Count},C_{h0r},C_{h0k},C_{h0N},C_{\rm gap},C_{\eta F},
 C_Y,C_{YL},C_{\rm pde},C_{ccK},C_F.
                                                               \tag{PN.40}
\end{gather*}
All are positive. For $c_{LL}$ this follows from
$1\leq c_{LL}^{-1}$ in the standard last-constant list; the other
nonfixed entries are at least one there or in the middle list.
The remaining fixed entries are
$C_{nW}=10$, $C_{Lw}=4455$, $C_{j0}=405$, $c_w=9/10$,
$C_{\eta F}=6$, $C_{YL}=9$ and $C_{ccK}=1$, with
$C_{\rm pde}\geq1$ supplied explicitly by standardness.
Thus the record's 23 positivity fields are also supplied.

The twenty-one families provide exactly 61 further named fields,
including fields that contain several quantified inequalities.
The remaining four fields are exactly the four clauses of the
explicit hypothesis $\operatorname{HostArith}(D.\mathrm{host})$:
\begin{enumerate}
\item for each $i$, an integral $\pi\in K$ with
      $\pi\overline\pi=g_i$;
\item $g_i\equiv1\pmod{QSqd}$ for each $i$;
\item the exact full splitting cardinalities above both auxiliary primes;
\item every supplied level prime has class $\mathrm{splitK}$.
\end{enumerate}
They become \texttt{g\_norm}, \texttt{g\_mod},
\texttt{pde\_split} and \texttt{level\_split}, respectively.
PN8's numerator conversion supplies the literal \texttt{e\_mu} field.
The ledger's field map checks every name in
\texttt{ParameterChoice'} exactly once:
\[
              23+23+61+4=111.                              \tag{PN.41}
\]
This is a concrete record, with the actual geometric data, embeddings,
weights and primes left in place. Since its constants were copied in
the order (PN.40), both required equalities hold exactly:
\[
                  \operatorname{parConsts}(P)=C.PK,
                  \qquad P.C_{nW}=C.P.C_{nW}.                \tag{PN.42}
\]
We have therefore proved the entire source conclusion
\[
 \operatorname{SzRow}\!\left(
 C,D,L\longmapsto
 \operatorname{HostArith}(D.\mathrm{host})\Longrightarrow
 \exists P:\operatorname{ParameterChoice}'(D,L),\
 \operatorname{parConsts}(P)=C.PK\ \land\ P.C_{nW}=C.P.C_{nW}
 \right).                                                    \tag{PN.43}
\]

\paragraph{The remaining existence boundary.}
No parameter family is retained after PN16. RF15 can now be applied to
this actual compatible $P$, with its already proved 54 \texttt{ParamsT}
fields; the numerical join with \texttt{SizesAt} is handled separately.
The hypothesis here is still an actual initial $D,L$ satisfying
\texttt{Val} and \texttt{HostArith}, including its rank windows,
host-prime windows, column identities, root-of-unity condition and
weight windows. Numerical estimates do not supply that host or those
columns. The four independent global theta/spectral, metric/dimension,
cohomological and final prime-window branches keep their earlier status.
